% GENERATED FILE. Edit canonical lessons, then run npm run generate:books.
% canonical-source-hash: fnv1a64:2ab1dab5d33f636c
% canonical-lessons: IT-C526-recensione, IT-C526-recensire, IT-C526-critico, IT-C526-divorare, IT-C526-rileggere, IT-R526-ripasso-519-526

\chapter{Leggere e giudicare}
\label{ch:it-leggere-e-giudicare}
\hlchaptermodality{526}

\begin{chapteropening}
\textbf{By the end of this chapter:} \emph{I can talk about reviews and the way I read.}

Complete the last lesson of chapter 526: I can talk about reviews and the way I read.
\end{chapteropening}

\section[la recensione]{la recensione — review}

\label{lesson:IT-C526-recensione}



\begin{quote}
Before the new one: say the Italian for bookmark, then the Italian for paperback.
\end{quote}

\subsection*{You'll want to know: la recensione}
\textbf{la recensione} — \textquotedblleft{}review\textquotedblright{}.

\begin{quote}
Ho letto una recensione molto positiva del film.
\end{quote}

\begin{quote}
\emph{I read a very good review of the film.}
\end{quote}

\emph{Scrivere una recensione} — to write a review.

\begin{cousinweb}
\emph{recensione} is a learned borrowing from Latin \emph{recēnsiō} \textquotedblleft{}a review, a counting\textquotedblright{} (Lewis \& Short), from \emph{cēnsēre} \textquotedblleft{}to assess\textquotedblright{}.
\end{cousinweb}

\begin{grammarlens}[title={What you've built}]
One more word for telling, explaining and describing at B1.
\end{grammarlens}

\subsection*{Your turn}
\begin{itemize}
\raggedright
  \item \emph{Say it:} \emph{la recensione}
  \item \emph{Say it:} \emph{la recensione}, once more
  \item \emph{Say it:} a sentence of your own with \emph{la recensione}
  \item \emph{Recall:} say the Italian for bookmark, then the Italian for paperback, then say \emph{la recensione} again
\end{itemize}

\begin{tcolorbox}[breakable,colback=teal!4,colframe=teal!35!black,title={Before you move on}]
What does la recensione mean? (\textquotedblleft{}Review\textquotedblright{}.) Say it once more.
\end{tcolorbox}
\section[recensire]{recensire — to review}

\label{lesson:IT-C526-recensire}



\begin{quote}
Before the new one: say the Italian for paperback, then the Italian for review.
\end{quote}

\subsection*{You'll want to know: recensire}
\textbf{recensire} — \textquotedblleft{}to review\textquotedblright{}.

\begin{quote}
Il giornale ha recensito il suo ultimo romanzo.
\end{quote}

\begin{quote}
\emph{The newspaper reviewed her latest novel.}
\end{quote}

An \emph{-isc-} verb: \emph{recensisco}.

\begin{cousinweb}
\emph{recensire} is a learned borrowing from Latin \emph{recēnsēre} \textquotedblleft{}to count over, to review\textquotedblright{} (Lewis \& Short), from \emph{cēnsēre} \textquotedblleft{}to assess\textquotedblright{}.
\end{cousinweb}

\begin{grammarlens}[title={What you've built}]
One more word for telling, explaining and describing at B1.
\end{grammarlens}

\subsection*{Your turn}
\begin{itemize}
\raggedright
  \item \emph{Say it:} \emph{recensire}
  \item \emph{Say it:} \emph{recensire}, once more
  \item \emph{Say it:} a sentence of your own with \emph{recensire}
  \item \emph{Recall:} say the Italian for paperback, then the Italian for review, then say \emph{recensire} again
\end{itemize}

\begin{tcolorbox}[breakable,colback=teal!4,colframe=teal!35!black,title={Before you move on}]
What does recensire mean? (\textquotedblleft{}To review\textquotedblright{}.) Say it once more.
\end{tcolorbox}
\section[il critico]{il critico — critic}

\label{lesson:IT-C526-critico}



\begin{quote}
Before the new one: say the Italian for review, then the Italian for to review.
\end{quote}

\subsection*{You'll want to know: il critico}
\textbf{il critico} — \textquotedblleft{}critic\textquotedblright{}.

\begin{quote}
Il critico ha parlato male dello spettacolo.
\end{quote}

\begin{quote}
\emph{The critic spoke badly of the show.}
\end{quote}

The plural is \emph{i critici}. Also an adjective: \emph{un momento critico}.

\begin{cousinweb}
\emph{critico} is a learned borrowing, through Latin \emph{criticus} (Lewis \& Short), from Greek \emph{kritikós} \textquotedblleft{}able to judge\textquotedblright{}, from \emph{krínein} \textquotedblleft{}to judge\textquotedblright{}.
\end{cousinweb}

\begin{grammarlens}[title={What you've built}]
One more word for telling, explaining and describing at B1.
\end{grammarlens}

\subsection*{Your turn}
\begin{itemize}
\raggedright
  \item \emph{Say it:} \emph{il critico}
  \item \emph{Say it:} \emph{il critico}, once more
  \item \emph{Say it:} a sentence of your own with \emph{il critico}
  \item \emph{Recall:} say the Italian for review, then the Italian for to review, then say \emph{il critico} again
\end{itemize}

\begin{tcolorbox}[breakable,colback=teal!4,colframe=teal!35!black,title={Before you move on}]
What does il critico mean? (\textquotedblleft{}Critic\textquotedblright{}.) Say it once more.
\end{tcolorbox}
\section[divorare]{divorare — to devour}

\label{lesson:IT-C526-divorare}



\begin{quote}
Before the new one: say the Italian for to review, then the Italian for critic.
\end{quote}

\subsection*{You'll want to know: divorare}
\textbf{divorare} — \textquotedblleft{}to devour\textquotedblright{}.

\begin{quote}
Ho divorato il libro in due giorni.
\end{quote}

\begin{quote}
\emph{I devoured the book in two days.}
\end{quote}

Of books: to read very fast.

\begin{cousinweb}
\emph{divorare} comes from Latin \emph{dēvorāre} \textquotedblleft{}to swallow down\textquotedblright{} (EtymDB), from \emph{vorāre} \textquotedblleft{}to swallow\textquotedblright{}.
\end{cousinweb}

\begin{grammarlens}[title={What you've built}]
One more word for telling, explaining and describing at B1.
\end{grammarlens}

\subsection*{Your turn}
\begin{itemize}
\raggedright
  \item \emph{Say it:} \emph{divorare}
  \item \emph{Say it:} \emph{divorare}, once more
  \item \emph{Say it:} a sentence of your own with \emph{divorare}
  \item \emph{Recall:} say the Italian for to review, then the Italian for critic, then say \emph{divorare} again
\end{itemize}

\begin{tcolorbox}[breakable,colback=teal!4,colframe=teal!35!black,title={Before you move on}]
What does divorare mean? (\textquotedblleft{}To devour\textquotedblright{}.) Say it once more.
\end{tcolorbox}
\section[rileggere]{rileggere — to reread}

\label{lesson:IT-C526-rileggere}



\begin{quote}
Before the new one: say the Italian for critic, then the Italian for to devour.
\end{quote}

\subsection*{You'll want to know: rileggere}
\textbf{rileggere} — \textquotedblleft{}to reread\textquotedblright{}.

\begin{quote}
Rileggo spesso i miei libri preferiti.
\end{quote}

\begin{quote}
\emph{I often reread my favourite books.}
\end{quote}

Like \emph{leggere}: past participle \emph{riletto}.

\begin{cousinweb}
\emph{rileggere} is built from \emph{ri-} \textquotedblleft{}again\textquotedblright{} and \emph{leggere} \textquotedblleft{}to read\textquotedblright{}, inherited from Latin \emph{legere} \textquotedblleft{}to gather, to read\textquotedblright{} (Lewis \& Short).
\end{cousinweb}

\begin{grammarlens}[title={What you've built}]
One more word for telling, explaining and describing at B1.
\end{grammarlens}

\subsection*{Your turn}
\begin{itemize}
\raggedright
  \item \emph{Say it:} \emph{rileggere}
  \item \emph{Say it:} \emph{rileggere}, once more
  \item \emph{Say it:} a sentence of your own with \emph{rileggere}
  \item \emph{Recall:} say the Italian for critic, then the Italian for to devour, then say \emph{rileggere} again
\end{itemize}

\begin{tcolorbox}[breakable,colback=teal!4,colframe=teal!35!black,title={Before you move on}]
What does rileggere mean? (\textquotedblleft{}To reread\textquotedblright{}.) Say it once more.
\end{tcolorbox}
\section[(dialogue)]{Ripasso dei capitoli 519–526}

\label{lesson:IT-R526-ripasso-519-526}



\begin{quote}
Say the Italian for stage and curtain.
\end{quote}

\subsection*{Your turn}
Say these aloud:

\begin{itemize}
\raggedright
  \item stage, curtain, stalls, gallery, applause
  \item to act, script, role, to make one's debut, director
  \item painting, paintbrush, canvas, easel, watercolour
  \item to sculpt, sculptor, marble, clay, to model
  \item exhibition, masterpiece, portrait, fresco, to restore
  \item novelist, chapter, to set, twist, ending
  \item bookseller, cover, to leaf through, bookmark, paperback
  \item review, to review, critic, to devour, to reread
\end{itemize}

\begin{tcolorbox}[breakable,colback=teal!4,colframe=teal!35!black,title={Before you move on}]
Stage? (\textbf{il palcoscenico}.) Curtain? (\textbf{il sipario}.) Stalls? (\textbf{la platea}.) Gallery? (\textbf{il loggione}.) Applause? (\textbf{l'applauso}.) To act? (\textbf{recitare}.) Script? (\textbf{il copione}.) Role? (\textbf{il ruolo}.) To make one's debut? (\textbf{debuttare}.) Director? (\textbf{il regista}.) Painting? (\textbf{la pittura}.) Paintbrush? (\textbf{il pennello}.) Canvas? (\textbf{la tela}.) Easel? (\textbf{il cavalletto}.) Watercolour? (\textbf{l'acquerello}.) To sculpt? (\textbf{scolpire}.) Sculptor? (\textbf{lo scultore}.) Marble? (\textbf{il marmo}.) Clay? (\textbf{l'argilla}.) To model? (\textbf{modellare}.) Exhibition? (\textbf{la mostra}.) Masterpiece? (\textbf{il capolavoro}.) Portrait? (\textbf{il ritratto}.) Fresco? (\textbf{l'affresco}.) To restore? (\textbf{restaurare}.) Novelist? (\textbf{il romanziere}.) Chapter? (\textbf{il capitolo}.) To set? (\textbf{ambientare}.) Twist? (\textbf{il colpo di scena}.) Ending? (\textbf{l'epilogo}.) Bookseller? (\textbf{il libraio}.) Cover? (\textbf{la copertina}.) To leaf through? (\textbf{sfogliare}.) Bookmark? (\textbf{il segnalibro}.) Paperback? (\textbf{il tascabile}.) Review? (\textbf{la recensione}.) To review? (\textbf{recensire}.) Critic? (\textbf{il critico}.) To devour? (\textbf{divorare}.) To reread? (\textbf{rileggere}.)
\end{tcolorbox}
